\documentclass[11pt,a4paper]{article}
\usepackage[top=1.7cm,bottom=1.5cm,left=1.5cm,right=1.5cm,headsep=8pt,footskip=16pt]{geometry}
\usepackage{enumitem}
\usepackage{amssymb}
\usepackage{array}
\usepackage[table]{xcolor}
\usepackage{fancyhdr}
\usepackage{lastpage}
\usepackage{multicol}
\usepackage{polyglossia}
\setdefaultlanguage[calendar=gregorian, numerals=maghrib]{arabic}
\setotherlanguage{english}
\IfFontExistsTF{Amiri}{%
  \newfontfamily\arabicfont[Script=Arabic, Scale=1.08]{Amiri}
}{%
  \newfontfamily\arabicfont[Script=Arabic, Scale=1.08, Path=../fonts/, Extension=.ttf,
    UprightFont=*-Regular, BoldFont=*-Bold, ItalicFont=*-Italic,
    BoldItalicFont=*-BoldItalic]{Amiri}
}
\IfFontExistsTF{AlBattar}{%
  \newfontfamily\khatt[Script=Arabic, Scale=1.15]{AlBattar}
}{%
  \newfontfamily\khatt[Script=Arabic, Scale=1.15, Path=../fonts/,
    BoldFont=ae_AlBattar.ttf]{ae_AlBattar.ttf}
}
\renewcommand{\headrulewidth}{0.8pt}
\renewcommand{\footrulewidth}{0.8pt}
\fancyhf{}
\rhead{\khatt المدرسة العليا للأساتذة ببوزريعة}
\lhead{\khatt ملحقة بومرداس}
\lfoot{\khatt الأستاذ: د. بن عبد الرزاق يوسف}
\cfoot{\thepage/\pageref{LastPage}}
\pagestyle{fancy}
\parindent 0pt
\setlength{\columnsep}{0.6cm}
\setlength{\columnseprule}{0.4pt}
\setlist[enumerate,1]{label=\textbf{\arabic*.}, leftmargin=*, itemsep=5pt, topsep=2pt}
\begin{document}

\linespread{1.08}\selectfont
\begin{center}
{\LARGE\khatt\textbf{الاختبار الرابع والعشرون في الذكاء الاصطناعي}}\\[4pt]
{\normalsize وضعيات إشكالية ــ أسئلة متعددة الاختيارات ــ الفصول من 1 إلى 6}
\end{center}
\vspace{-2pt}\rule{\textwidth}{0.6pt}\\[2pt]
\textbf{التعليمة:} يتضمن الاختبار خمس وضعيات إشكالية. اقرأ كل وضعية بعناية، ثم اختر التحليل أو التصرف الأنسب
من بين الاختيارات الخمسة؛ لكل وضعية إجابة صحيحة واحدة فقط.
اكتب \textbf{حرف} الإجابة الصحيحة في جدول الإجابات في آخر الاختبار.
لكل وضعية 4 نقاط (المجموع 20 نقطة).
\vspace{4pt}

\noindent
{\large\khatt\textbf{الوضعية 1}}\hfill{\small(الفصل 2)}\par\nopagebreak\vspace{3pt}
\setlength{\fboxsep}{7pt}\colorbox{gray!10}{\parbox{\dimexpr\linewidth-14pt}{لاحظ أستاذ أن أداة ذكاء اصطناعي تجيب عن أسئلة البلاغة الغربية بالإنجليزية بدقة وتفصيل، بينما تأتي إجاباتها عن مسائل البلاغة العربية بالعربية أقل دقة وأكثر عمومية. ويريد أن يفهم السبب ليقرر طريقة استعمالها مع طلبته. وطلبته في السنة الثالثة ليسانس يعتمدون كثيرًا على هذه الأداة في بحوثهم. أي تفسير هو الأقرب إلى طريقة تعلّم النماذج؟}}\par\nopagebreak
\begin{itemize}[leftmargin=1.0cm, labelsep=6pt, itemsep=3pt, topsep=4pt, parsep=0pt]
  \item[\textbf{أ)}] الحل أن يطرح أسئلة البلاغة العربية بالإنجليزية، لأن الإجابة ستأتي حينئذ دقيقة ومفصّلة بالضرورة، فالموضوع لا يؤثر في الأداء، واللغة المستعملة في السؤال هي العامل الوحيد.
  \item[\textbf{ب)}] كمية البيانات تؤثر في الأداء، ولذلك تُعتمد إجابات البلاغة الغربية دون تحقق، أما إجابات البلاغة العربية فلا يُتحقق منها أيضًا، لأن الأستاذ متخصص ويكتشف أخطاءها عند القراءة.
  \item[\textbf{ج)}] البلاغة العربية أصعب بطبيعتها من البلاغة الغربية، ولذلك يستحيل على أي نموذج إتقانها مهما كانت بياناته، فيتجنب الأستاذ أسئلة البلاغة العربية تمامًا مع الأداة ومع طلبته.
  \item[\textbf{د)}] كمية البيانات وجودتها في كل لغة وموضوع تؤثر في أداء النموذج، والمحتوى الإنجليزي أوفر في بيانات التدريب عادة؛ فيُتحقق بعناية أكبر من الإجابات في المسائل العربية المتخصصة.
  \item[\textbf{هـ)}] النموذج يفضّل الإنجليزية عمدًا لأن مطوّريه برمجوه يدويًا على تقليل جودة الإجابات العربية، ولا علاقة لذلك بكمية البيانات أو جودتها في كل لغة، فلا يمكن تحسين النتائج بأي طريقة.
\end{itemize}
\par\vspace{6pt}
\noindent
{\large\khatt\textbf{الوضعية 2}}\hfill{\small(الفصل 3)}\par\nopagebreak\vspace{3pt}
\setlength{\fboxsep}{7pt}\colorbox{gray!10}{\parbox{\dimexpr\linewidth-14pt}{يحضّر طلبة مناظرة حول «فرض الفصحى في كل مواد التعليم»، وكل فريق يطلب من الأداة حججًا تؤيد موقفه فقط. لاحظ الأستاذ أن الفرق لا تتوقع حجج خصومها فتنهار في الردود، وأن بعض الحجج المولَّدة مدعومة بإحصاءات غير موثقة. أي استعمال للأداة أنسب لتحضير المناظرة؟}}\par\nopagebreak
\begin{itemize}[leftmargin=1.0cm, labelsep=6pt, itemsep=3pt, topsep=4pt, parsep=0pt]
  \item[\textbf{أ)}] يطلب كل فريق أقوى الحجج المؤيدة والمعارضة، وتوقع الاعتراضات واقتراح ردود، ثم يتحقق من كل إحصاء في مصادر موثوقة، ويصوغ حججه بنفسه، لأن المناظرة تقيس تفكيره لا مخرجات الأداة.
  \item[\textbf{ب)}] يطلب كل فريق من الأداة أن تؤدي المناظرة كاملة نيابة عنه، ثم يحفظ النص المولَّد ويلقيه، لأن ذلك يضمن حججًا وردودًا متكاملة، ولا حاجة إلى التحقق من الإحصاءات.
  \item[\textbf{ج)}] يطلب كل فريق الحجج المؤيدة والمعارضة، ثم يعتمد الإحصاءات كما وردت دون تحقق، لأن الأداة لا تقدّم في سياق المناظرات إلا أرقامًا موثقة من دراسات منشورة ومحكّمة، فتنال الفرق علامات أعلى في الحجاج.
  \item[\textbf{د)}] يطلب كل فريق حججًا مؤيدة أكثر عددًا وقوة، لأن كثرة الحجج المؤيدة تغني عن توقع حجج الخصوم، والإحصاءات غير الموثقة تزيد قوة الإقناع أمام لجنة التحكيم في المناظرة، وتبهر الجمهور والمحكّمين معًا.
  \item[\textbf{هـ)}] يمنع الأستاذ أي استعمال للأداة في التحضير، لأن الأدوات لا تستطيع إنتاج حجج معارضة لموقف المستخدم، فهي تؤيد دائمًا ما يطلبه، ولا تصلح لتوقع الاعتراضات أبدًا.
\end{itemize}
\par\vspace{6pt}
\noindent
{\large\khatt\textbf{الوضعية 3}}\hfill{\small(الفصل 4)}\par\nopagebreak\vspace{3pt}
\setlength{\fboxsep}{7pt}\colorbox{gray!10}{\parbox{\dimexpr\linewidth-14pt}{يدرّس أستاذ مئتي طالب، فاستعمل أداة ذكاء اصطناعي لتوليد ملاحظات «شخصية» على بحوثهم دون أن يقرأ أغلبها. لاحظ الطلبة أن الملاحظات متشابهة، وأن بعضها يمدح عناصر غير موجودة في بحوثهم، فاحتجّوا لدى رئيس القسم. ويريد رئيس القسم أن يحدد موقفًا واضحًا من هذه الممارسة لكل أساتذة القسم. ما التقدير الأدق؟}}\par\nopagebreak
\begin{itemize}[leftmargin=1.0cm, labelsep=6pt, itemsep=3pt, topsep=4pt, parsep=0pt]
  \item[\textbf{أ)}] لا إشكال، لأن عدد الطلبة كبير، والأداة اطلعت على البحوث وقدّمت ملاحظات مناسبة، والتشابه بين الملاحظات طبيعي لأن البحوث متشابهة في أغلبها، فلا يلزم الأستاذ قراءتها.
  \item[\textbf{ب)}] الإشكال في عدم الإفصاح فقط، فإذا أخبر الطلبة أن الملاحظات مولَّدة صار التقويم سليمًا، حتى لو لم يقرأ البحوث، لأن الإفصاح ينقل المسؤولية من الأستاذ إلى الأداة.
  \item[\textbf{ج)}] تقديم ملاحظات مولَّدة على أنها تقويم شخصي دون قراءة البحوث يخل بمسؤولية الأستاذ ويضلل الطلبة؛ ويمكنه الاستعانة بالأداة في الصياغة بعد قراءته وتحديده الملاحظات، مع الإفصاح.
  \item[\textbf{د)}] الإشكال الوحيد في المدح غير المطابق، فيكفي أن يحذف عبارات المدح من الملاحظات المولَّدة ويرسل الباقي، دون أن يقرأ البحوث، لأن الملاحظات النقدية المولَّدة دقيقة دائمًا.
  \item[\textbf{هـ)}] يُمنع الأستاذ من أي استعمال للأدوات في التقويم، حتى في تحسين صياغة ملاحظات كتبها بنفسه بعد قراءة البحوث، لأن مجرد استعمال الأداة يُفقد الملاحظة طابعها الشخصي.
\end{itemize}
\par\vspace{6pt}
\noindent
{\large\khatt\textbf{الوضعية 4}}\hfill{\small(الفصل 5)}\par\nopagebreak\vspace{3pt}
\setlength{\fboxsep}{7pt}\colorbox{gray!10}{\parbox{\dimexpr\linewidth-14pt}{كتب طالب في بحثه رقمًا محددًا لـ«عدد الناطقين بالعربية» قدّمته أداة ذكاء اصطناعي دون تاريخ أو مصدر. ولاحظ أستاذه أن مصادر مختلفة تذكر أرقامًا متباينة كثيرًا، وأن بعضها يحسب الناطقين بها لغةً أولى فقط، وبعضها يضيف من يتقنها لغةً ثانية. كيف يعالج الطالب هذه المعلومة؟}}\par\nopagebreak
\begin{itemize}[leftmargin=1.0cm, labelsep=6pt, itemsep=3pt, topsep=4pt, parsep=0pt]
  \item[\textbf{أ)}] يسأل الأداة عن مصدر الرقم، فإن ذكرت منظمة دولية أحال إليها مباشرة دون الرجوع إلى تقريرها، لأن ذكر اسم المنظمة يكفي للتوثيق في البحوث الجامعية.
  \item[\textbf{ب)}] يرجع إلى مصادر إحصائية موثوقة، ويذكر الرقم مع مصدره وسنته وتعريفه، لغةً أولى أو ثانية، ويشير إلى تباين التقديرات، لأن هذه المعطيات تتغير وتختلف باختلاف منهجية الحساب.
  \item[\textbf{ج)}] يحذف الرقم ويكتب «عدد كبير جدًا»، لأن الأرقام الإحصائية متباينة دائمًا، فلا يمكن ذكر أي رقم دقيق في بحث جامعي، حتى لو أُرفق بمصدره وسنته وتعريفه.
  \item[\textbf{د)}] يختار أكبر رقم وجده في المصادر، لأنه يعكس مكانة العربية، ويذكره دون سنة أو تعريف، لأن الأرقام الكبيرة لا تحتاج إلى توضيح منهجية حسابها في البحوث الجامعية، فيخدم بذلك قضية البحث.
  \item[\textbf{هـ)}] يُبقي الرقم الذي قدّمته الأداة، لأنها جمعت كل المصادر وحسبت المتوسط، وذكر السنة والتعريف تفاصيل ثانوية لا يحتاجها القارئ في بحث عن اللغة العربية وانتشارها.
\end{itemize}
\par\vspace{6pt}
\noindent
{\large\khatt\textbf{الوضعية 5}}\hfill{\small(الفصل 6)}\par\nopagebreak\vspace{3pt}
\setlength{\fboxsep}{7pt}\colorbox{gray!10}{\parbox{\dimexpr\linewidth-14pt}{استنسخت أستاذة صوتها بأداة ذكاء اصطناعي لتسجيل دروس صوتية لتلاميذها خلال عطلتها المرضية، دون أن تخبرهم بأن الصوت مولَّد. وبعد أيام وصل الأولياءَ تسجيل بصوتها يطلب مبلغًا ماليًا «لنشاط مدرسي»، لم تسجله هي. والإدارة تنتظر منها توضيحًا عاجلًا، وبعض الأولياء قلقون ويطالبون بتفسير. ما الدرس الأهم والتصرف الأنسب؟}}\par\nopagebreak
\begin{itemize}[leftmargin=1.0cm, labelsep=6pt, itemsep=3pt, topsep=4pt, parsep=0pt]
  \item[\textbf{أ)}] تحذف دروسها الصوتية كلها وتعتذر عن استعمال الذكاء الاصطناعي، لأن المسؤولية عن التسجيل المزيف تقع عليها وحدها، ولا حاجة إلى إبلاغ الإدارة أو الأولياء أو مراجعة الحماية.
  \item[\textbf{ب)}] تُبلغ الأولياء بأن التسجيل مزيف، لكنها تواصل الاستنساخ دون إفصاح أو مراجعة لحماية حسابها، لأن التزييف حادث فردي لا يتكرر، والإفصاح عن طبيعة الصوت يُضعف ثقة التلاميذ.
  \item[\textbf{ج)}] تتجاهل التسجيل المزيف، لأن الأولياء سيكتشفون وحدهم أنه ليس صوتها، فلا داعي لإبلاغ الإدارة، وتواصل استعمال الصوت المستنسخ دون إفصاح، لأن ما حدث لا علاقة له بدروسها.
  \item[\textbf{د)}] تطلب من الأولياء دفع المبلغ مؤقتًا حتى تتأكد الإدارة من مصدر التسجيل، لأن الطلب قد يكون صحيحًا صادرًا عن زميلة استعملت صوتها لغرض تربوي مشروع.
  \item[\textbf{هـ)}] الصوت المستنسخ قد يُساء استعماله، فتُبلغ الإدارة والأولياء فورًا بأن التسجيل مزيف، وتراجع حماية حسابها وملفاتها، وتلتزم بالإفصاح عن طبيعة الصوت وبقناة رسمية للمراسلات المالية.
\end{itemize}
\par\vspace{6pt}
\begin{center}
{\khatt\textbf{جدول الإجابات}}\\[3pt]
\renewcommand{\arraystretch}{1.6}
\begin{tabular}{|>{\centering\arraybackslash}m{1.9cm}|*{5}{>{\centering\arraybackslash}m{1.2cm}|}}
\hline
\textbf{رقم السؤال} & \textbf{1} & \textbf{2} & \textbf{3} & \textbf{4} & \textbf{5} \\ \hline
\textbf{الإجابة} &  &  &  &  &  \\[8pt] \hline
\end{tabular}

\vspace{8pt}
\renewcommand{\arraystretch}{2.2}
\begin{tabular}{|>{\centering\arraybackslash}m{7.5cm}|>{\centering\arraybackslash}m{3.2cm}|>{\centering\arraybackslash}m{3.2cm}|}
\hline
\textbf{الاسم واللقب} & \textbf{الفوج} & \textbf{العلامة} \\ \hline
 & & \hfill /20 \\ \hline
\end{tabular}
\end{center}
\end{document}
